% HIKET -- PROPOSAL: replacing the independent-observation likelihood with a
% correlated-error (compound symmetry) likelihood.
%
% STATUS: agreed plan, not yet implemented. Written 2026-08-17, REWRITTEN 2026-08-18
%         after the auxiliary-sigma-prior run landed and the design was worked out.
% Scope: this document describes ONE change to ONE function. It deliberately spells
% out every step, including the ones that are obvious to someone who works with
% mixed models daily, because the argument has to survive review by people who do not.
\documentclass[11pt]{article}
\usepackage[utf8]{inputenc}\usepackage[T1]{fontenc}
\usepackage[margin=2.4cm]{geometry}
\usepackage{graphicx}\usepackage{amsmath}\usepackage{amssymb}\usepackage{xcolor}
\usepackage{booktabs}\usepackage{longtable}
\usepackage{float}\usepackage[hidelinks]{hyperref}
\usepackage{caption}\captionsetup{font=small,labelfont=bf}
\usepackage{enumitem}
\usepackage{framed}

% Dependency-light callout boxes: framed + a bold lead line. Deliberately avoids
% tcolorbox, which failed to build on the authoring machine (pgfkeys '/tcb/not after').
\newenvironment{keybox}[1][]{%
  \begin{leftbar}\noindent\textbf{\small #1}\par\vspace{0.3em}\small}{\end{leftbar}}
\newenvironment{warnbox}[1][]{%
  \begin{leftbar}\noindent\textcolor{red!60!black}{\textbf{\small #1}}\par\vspace{0.3em}\small}{\end{leftbar}}

\title{Replacing the independent-observation likelihood\\
\large A proposal for the HIKET calibration}
\author{}
\date{Revised 18 August 2026 --- \emph{agreed plan, not yet implemented}}

\begin{document}
\maketitle

\begin{keybox}[In one paragraph]
The calibration currently assumes that every soil-carbon measurement is an
independent piece of evidence. It is not. The same plot measured in three campaigns
gives three strongly related numbers; plots near each other give related numbers
too; and every plot in one campaign shares whatever that campaign's method did.
Counting 1205 measurements as 1205 independent facts overstates the evidence by
roughly a factor of five for the quantity that controls the headline results. This
document sets out the agreed plan: three shared-offset terms, all with their sizes
fixed from outside the calibration, replacing the independence assumption. It gives
the numbers, the traps found while deriving them, what the change is expected to
achieve, and --- equally important --- what it cannot.
\end{keybox}

\section{What the calibration does now}

For plot $i$ in campaign $j$ the model returns a predicted carbon stock
$f_{ij}(\theta)$, and the comparison with the observation is made on the log scale:
\[
\log y_{ij} \;=\; \log f_{ij}(\theta) \;+\; e_{ij},
\qquad e_{ij}\sim N(0,\sigma^2), \quad \sigma = 0.800 .
\]
The log-likelihood is the plain sum of the individual log-densities. That sum
\emph{is} the assumption that every $e_{ij}$ is independent of every other. The
assumption has never been stated as a hypothesis anywhere in the code or the
manuscript; it arrived as a default.

\section{Why that assumption is wrong here}

Three different groupings make observations share information. They are easy to
conflate, so we name them explicitly and use one notation throughout: a shared
offset is $u$ with a superscript naming what it groups.

\begin{center}
\begin{tabular}{llrl}
\toprule
term & groups & members & meaning \\
\midrule
$u^{R}_{r}$ & plots in a region & $\sim$150 & same climate, parent material, forest history \\
$u^{P}_{i}$ & visits to one plot & 2.6 & same soil, missing edaphic control \\
$u^{C}_{j}$ & all plots in one campaign & $\sim$400 & shared method, crew, laboratory \\
\bottomrule
\end{tabular}
\end{center}

\begin{keybox}[The counter-intuitive part]
A \emph{small} correlation shared by \emph{many} observations does more damage to a
national mean than a large one shared by few. The plot term has a correlation of
0.5--0.6 but only 2.6 members. The regional term has a correlation of just
\textbf{0.080} --- negligible-sounding --- but $\sim$150 plots share it, and
$1+149\times0.08 \approx 13$. Group size multiplies the effect. This is why the
regional term matters more than the plot term, which is the opposite of what the
raw correlations suggest.
\end{keybox}

\subsection{Measured, from the observations rather than from residuals}

\begin{keybox}[Design principle]
Residuals may \textbf{check} a design value. They may never \textbf{set} one.
Setting the likelihood from the residuals of a fit produced under the assumption we
are replacing is circular. Every number below comes from the observations
themselves, with campaign levels removed first so nothing can be a trend artefact.
\end{keybox}

\paragraph{The plot term.} Under compound symmetry the covariance between any two
visits to the same plot \emph{is} $\tau_P^2$, so it can be read directly off a
campaign pair with no variance decomposition and no distributional assumption:

\begin{center}
\begin{tabular}{lrrr}
\toprule
pair & $n$ & correlation & $\tau_P$ \\
\midrule
1985--2006 & 384 & 0.805 & 0.355 \\
1985--2024 & 342 & 0.762 & 0.363 \\
\textbf{2006--2024} & \textbf{442} & \textbf{0.820} & \textbf{0.396} \\
\bottomrule
\end{tabular}
\end{center}

We adopt $\tau_P = 0.396$ from the 2006--2024 pair: the largest sample, the widest
of the three, and \textbf{the only pair untouched by the 1985 problems} --- no
imputed litter layer, no 1986--1995 dating spread. The 1985 campaign therefore
enters the likelihood at full weight but sets none of the likelihood's parameters.

\paragraph{The regional term.} Splitting the plots into eight equal-count latitude
bands (ETRS northing), after removing campaign levels, the band means run
monotonically south to north:
\[
+0.18,\;+0.16,\;+0.06,\;+0.01,\;-0.03,\;-0.16,\;-0.12,\;-0.10
\]
a gradient of about 32\%. Between-band $\mathrm{sd} = \mathbf{0.117}$, between-plot
within band $=0.397$, giving a regional correlation of 0.080. We adopt
$\tau_R = 0.117$.

\paragraph{The campaign term.} Not measurable from this data at all --- see
Section~\ref{sec:campaign}.

\subsection{There is no field replication to check any of this}

The Biosoil workbook carries a \texttt{REP} column, but only 63 rows of 3028 (about
16 plots) have \texttt{REP}~$=2$, all in Biosoil, none in 1985 or 2024. Of the three
carbon columns, two return $\mathrm{sd}(\log\text{ratio}) = 0.000$ and $0.003$
between the pairs --- copied values, not independent re-measurements. Only one shows
genuine variation, at about 12\% per sample, and that reads as a laboratory
duplicate rather than a second field sample.

\begin{warnbox}[Consequence]
We cannot empirically separate ``we sampled a different spot'' from ``this plot
really is different''. Random relocation between campaigns is harmless for the
design --- it is genuinely independent between visits, so it already sits inside
$\sigma_e$ --- but a \emph{systematic} component of relocation or protocol drift is
a campaign-level offset, and nothing in the data can measure it.
\end{warnbox}

\section{Why simply assuming a larger error does not work}

Inflating $\sigma$ multiplies every direction of the posterior by $\sqrt{k}$.
Reaching what is needed would require $\sigma \approx 1.8$ against a measured
residual spread of 0.71, which would destroy the self-consistency that justifies
$\sigma$ in the first place. \textbf{Correlation widens the level at fixed total
variance; nothing else available does.} That is the whole argument for this
instrument, and it is why the variance is \emph{split} rather than \emph{added}:
\[
\sigma^2_{\text{total}} \;=\; \tau_R^2 + \tau_P^2 + \tau_C^2 + \sigma_e^2 ,
\qquad \sigma_{\text{total}} = 0.800 \text{ (unchanged)} .
\]

\begin{warnbox}[A trap found while deriving this --- do not repeat it]
The observations show an intraclass correlation of \textbf{0.799}. Transplanting
that \emph{ratio} into a total of 0.800 gives $\tau = 0.715$ and $\sigma_e = 0.359$
--- \emph{below} the models' own within-plot error of 0.44--0.48. That would assert
the trend is measured more precisely than any model can reproduce it: the same
overconfidence we are removing from the level, recreated in the trend. The 0.799 is
a ratio inside the observations' own spread of 0.422; only its \emph{absolute} part
transfers. Always carry absolute standard deviations across, never ratios.
\end{warnbox}

\section{The plan}

\subsection{The model}

\[
\log y_{ij} \;=\; \log f_{ij}(\theta) \;+\; u^{R}_{r(i)} \;+\; u^{P}_{i} \;+\;
u^{C}_{j} \;+\; e_{ij}
\]
with $u^{R}\sim N(0,\tau_R^2)$, $u^{P}\sim N(0,\tau_P^2)$,
$u^{C}\sim N(0,\tau_C^2)$, $e\sim N(0,\sigma_e^2)$.

\subsection{Nothing new is estimated}

None of the offsets is ever given a value. They are \textbf{marginalised} ---
integrated out, averaged over all the values they might take --- leaving behind only
their effect on the covariance:
\[
\Sigma \;=\; \sigma_e^2 I \;+\; \tau_P^2 Z_P Z_P' \;+\; \tau_R^2 Z_R Z_R'
\;+\; \tau_C^2 Z_C Z_C' .
\]
This distinction is the crux of the whole proposal and is worth stating plainly:

\begin{center}
\begin{tabular}{lll}
\toprule
 & what happens to the offset & consequence \\
\midrule
ignore it & fixed at exactly zero & asserts perfect comparability \\
\textbf{fit it} & data choose the best-fitting value & \textbf{eats the trend} ($\delta=-0.137$) \\
\textbf{marginalise it} & never given any value & widens the right direction \\
\bottomrule
\end{tabular}
\end{center}

The standing rule that ``a campaign-specific bias is forbidden'' applies to the
\emph{fitted} version, which is unidentifiable and absorbs the signal. The
marginalised version adds no parameter and cannot absorb anything, because there is
nothing to absorb \emph{with}.

\subsection{The numbers}

\begin{center}
\begin{tabular}{llll}
\toprule
term & value & source & status \\
\midrule
$\tau_R$   & 0.117 & latitude-band means, campaign-centred & measured, fixed \\
$\tau_P$   & 0.396 & 2006--2024 pair covariance & measured, fixed \\
$\tau_C$   & 0.06 (1985), 0.03 (2006, 2024) & prescribed, see below & \textbf{assumed} \\
$\sigma_e$ & 0.685 & remainder of the fixed total & derived \\
\bottomrule
\end{tabular}
\end{center}

\begin{keybox}[The one-sided check]
$\sigma_e = 0.685$ sits comfortably above the models' measured within-plot error of
0.44--0.48, so the trend is not over-claimed. This is a residual-based
\emph{check}, not a residual-based \emph{setting} --- the distinction the design
principle allows.
\end{keybox}

\subsection{Setting $\tau_C$}
\label{sec:campaign}

$\tau_C$ cannot be measured; it is chosen. That is not a weakness relative to the
status quo, because the status quo also chooses a value --- zero --- and zero is the
least defensible number available, asserting that three campaigns spanning 39 years,
two laboratory methods and three field crews are perfectly inter-comparable.

The choice rests on two arguments that happen to agree. First, the largest
campaign-level defect we have actually identified and corrected --- the 1985
litter-layer definition --- was worth 3.37--4.27\,Mg\,C\,ha$^{-1}$, i.e.\
\textbf{4.8--6.0\% of stock}, and at least three further mechanisms are documented
but unquantified (LOI-derived mineral carbon in 1985 versus dry combustion later,
relocated subplots, the 1986--1995 dating spread). The residual should be
\emph{at least} the size of the one we happened to measure. Second, 1985 deserves
roughly twice the others: different laboratory method, relocated subplots, an
imputed litter layer, a ten-year dating spread, and no external validation, whereas
the 2006 and 2024 processing reproduces LUKE's official national stocks exactly.
A 3\% base doubled for 1985 gives 6\%, which is also the top of the measured
litter-layer range. \textbf{This is a stated assumption, not a measurement, and
should be written as one.}

\section{What the change is expected to achieve}

\subsection{Effect on the evidence}

On the calibration target (1205 observations, 456 plots; 310 with three visits,
129 with two, 17 with one):

\begin{center}
\begin{tabular}{lrrr}
\toprule
structure & SE of the national mean & inflation & effective $n$ \\
\midrule
independent (now) & 0.0231 & 1.00 & 1205 \\
$+\,u^{P}$        & 0.0274 & 1.19 & 855 \\
$+\,u^{P}+u^{R}$  & 0.0496 & \textbf{2.15} & \textbf{260} \\
\bottomrule
\end{tabular}
\end{center}

\begin{warnbox}[The plot term alone is not worth doing]
$u^{P}$ by itself reduces precision by a factor 1.4, against a likelihood that is
already 10--15 times more informative than the prior for $\sigma_{\text{input}}$. It
moves the estimates by 1--2\%. \textbf{The regional term is where the effect is}, and
the two must go in together.
\end{warnbox}

\subsection{Expected landing}

Marginal, local approximations --- direction is predictable, magnitude is not; an
estimate of exactly this kind gave an effective sample size of 1 in 1001 when used
to predict a prior change. Treat as order of magnitude.

\begin{center}
\begin{tabular}{lrr}
\toprule
 & now & expected \\
\midrule
posterior sd of $\log\sigma_{\text{input}}$ & 0.062--0.076 & roughly double \\
$\sigma_{\text{input}}$ & 1.69--2.10 & 1.51--1.80 \\
intrinsic MRT (Yasso07/15/20) & 21.7 / 25.1 / 21.4 & $\sim$25.4 / 28.5 / 24.8 \\
effective litter flux & 4.25--5.28 & 3.8--4.6 \\
\bottomrule
\end{tabular}
\end{center}

The flux row is the checkable one: all six models would fall inside the Nordic
Class~I window of 4.62\,tC\,ha$^{-1}$\,yr$^{-1}$ for the first time.

\begin{warnbox}[No prediction is offered for $\sigma_{\text{init}}$]
The arithmetic above assumes the likelihood loosens uniformly in every direction.
That is apt for a level parameter such as $\sigma_{\text{input}}$ and \emph{wrong}
for $\sigma_{\text{init}}$, whose information lives in trajectory shape --- and the
shared offsets cancel on differencing, leaving within-plot contrasts untouched.
$\sigma_{\text{init}}$ should move only as a knock-on through its correlation with
$\sigma_{\text{input}}$, measured at $-0.29$ to $-0.45$. An earlier reading of the
marginal prior-to-posterior widths concluded the opposite and was wrong.
\end{warnbox}

\section{What it will not achieve}

\subsection{The trend is a separate problem, and the short interval is the fragile one}

The correlated likelihood widens the \emph{level}. It does nothing for the
\emph{trend}, which is governed by campaign comparability. Writing the implied
uncertainty on an observed rate as $\tau_C\sqrt{2}\,\bar{C}/\Delta t$:

\begin{center}
\begin{tabular}{lrrr}
\toprule
 & signal & \% of stock & $\tau_C$ that dissolves it (1$\sigma$) \\
\midrule
1985$\rightarrow$2024 & 9.1\,Mg\,ha$^{-1}$ over 35.1\,yr & 12.8 & \textbf{9.1\%} \\
2006$\rightarrow$2024 & 2.1\,Mg\,ha$^{-1}$ over 18.0\,yr & \textbf{3.0} & \textbf{2.1\%} \\
\bottomrule
\end{tabular}
\end{center}

\begin{center}
\begin{tabular}{lrr}
\toprule
$\tau_C$ scheme (1985 / others) & 1985--2024 ($+0.259$) & 2006--2024 ($+0.117$) \\
\midrule
0 (current, implicit) & exact & exact \\
4\% / 2\%  & $\pm0.090$ \;(2.87$\sigma$) & $\pm0.111$ \;(1.05$\sigma$) \\
\textbf{6\% / 3\% (adopted)} & $\pm0.135$ \;(1.91$\sigma$) & $\pm0.167$ \;(0.70$\sigma$) \\
10\% / 5\% & $\pm0.226$ \;(1.15$\sigma$) & $\pm0.278$ \;(0.42$\sigma$) \\
\bottomrule
\end{tabular}
\end{center}

\begin{keybox}[The main result of this analysis, and it is not about the likelihood]
The intuition has been that 1985 is the campaign to distrust, so the long
comparison is the weak one. The arithmetic says the reverse. \textbf{1985--2024 is
robust} --- 35 years is long enough for the signal to outgrow any plausible bias,
surviving to $\tau_C \approx 9\%$. \textbf{2006--2024 is not resolvable} at any
$\tau_C$ above about 2\%, even though it uses our two best campaigns: same protocol,
both reproducing LUKE's official stocks exactly. An 18-year change of 3\% of stock
simply cannot survive a campaign offset of comparable size. No amount of trusting
those two campaigns repairs this; the interval is too short relative to the signal.
\end{keybox}

\begin{warnbox}[A recorded result that must be restated]
The finding that ``2006$\rightarrow$2024 is a source in all six models'' has died,
with four models now bracketing the observed $+0.117$. That is much weaker than it
reads: $+0.117$ itself is not resolvable against any campaign bias above $\sim$2\%.
``The models bracket the observation'' is not evidence when the observation carries
an unquantified systematic of its own size.
\end{warnbox}

This table is analytic --- arithmetic on the campaign means --- so all three
$\tau_C$ schemes are reported at no computational cost, whatever is run.

\subsection{Other things it does not fix}

The plot term is not a substitute for edaphic predictors; no model in the ensemble
has any, which is why $R^2 \approx 0.02$ and why a persistent per-plot offset is the
right shape for the missing information. And nothing here addresses the possibility
that the models are structurally wrong about early accumulation --- the misfit
concentrated in 1985--2006 remains whatever the likelihood says.

\section{Implementation}

\subsection{It is computationally trivial, for one specific reason}

Every variance component is \textbf{fixed}. $\Sigma$ therefore does not depend on
$\theta$ and does not change between MCMC draws. Factorise it \emph{once} at setup
--- a $1205\times1205$ Cholesky, well under a second, about 12\,MB stored --- and
each likelihood evaluation is a triangular solve for $r'\Sigma^{-1}r$ with
$r = \log y - \log f(\theta)$. That is of order 1\,ms against a forward model
costing 227\,ms: under half a percent of overhead. $\log|\Sigma|$ is computed once.

Woodbury and Sherman--Morrison identities are \emph{not} required. They would earn
their place only if the variance components were being estimated, or if $n$ were far
larger.

\begin{warnbox}[The one structural change]
The likelihood is currently computed plot-by-plot under \texttt{mclapply} and
summed. The plot term preserves that decomposition and the regional term preserves
it with regions as blocks, but \textbf{the campaign term couples every plot and
every region simultaneously}, so the per-plot sum is replaced by a single global
solve. This is a change to how the likelihood is assembled, not a local edit.
\end{warnbox}

\subsection{Build order --- these are local unit tests, not calibrations}

Each check is a single likelihood evaluation at a fixed parameter vector, compared
against the previous stage. Seconds on a laptop; no MCMC.

\begin{enumerate}[nosep]
\item $\tau_P = \tau_R = \tau_C = 0$ must reproduce the current log-likelihood
      \textbf{exactly}, to floating point. This is the only test that separates a
      correct implementation from a plausible-looking one.
\item Add $u^{P}$; setting $\tau_P=0$ must return to step 1.
\item Add $u^{R}$; setting $\tau_R=0$ must return to step 2.
\item Add $u^{C}$; setting $\tau_C=0$ must return to step 3.
\end{enumerate}

\subsection{Run configuration}

\begin{center}
\begin{tabular}{lll}
\toprule
switch & value & note \\
\midrule
correlation & $u^{R}+u^{P}+u^{C}$ & values in the table above \\
\texttt{HIKET\_SIGMA\_TOTAL} & 0.800 & unchanged; split, not added \\
\texttt{HIKET\_LIK\_DF} & \textbf{unset} & Student-$t$ dropped, see below \\
\texttt{HIKET\_SIGMA\_1985\_INFL} & \textbf{1.0} & replaced by $\tau_C$(1985)=0.06 \\
chains / iterations & 5 $\times$ 50000 & unchanged \\
\bottomrule
\end{tabular}
\end{center}

\textbf{One launch, six jobs, all six models.} No staged runs: the intermediate
stages are unit tests, not calibrations.

\paragraph{Why Student-$t$ is dropped.} The tails were adopted because measured
kurtosis is 6.9--7.1 against a Gaussian 3. But pooling observations from plots with
different persistent offsets produces heavy tails \emph{mechanically} --- a mixture
of Gaussians with different means is leptokurtic. Tails and shared offsets are
competing explanations for the same measurement, and fitting both risks paying twice
for one structure. Drop them, then measure kurtosis \emph{conditional on} the fitted
offsets: if it remains near 7 the tails are independently justified and return; if
it falls toward 3 they were a proxy for exactly this. Note that $\sigma$ means the
same thing either way, since the current setup already maps it to the standard
deviation.

\paragraph{The 1985 inflation is replaced, not removed.} 1985 keeps its extra
weight --- $\tau_C = 0.06$ against 0.03 for the others, the factor of two argued for
above. What is switched off is the old \emph{device}. \texttt{SIGMA\_1985\_INFL}~=~2.0
inflated the \emph{independent, per-observation} noise for 1985, reaching the campaign
level only because $n$ is finite (as $1/\sqrt{n}$) and degrading every plot-level
comparison in 1985 on the way. $\tau_C$ is a \emph{shared} level offset: it displaces
the campaign level and does nothing else, whatever $n$ is. Since the documented defect
is a level and not scatter --- 1985's noise level is normal, subsoil repeatability 0.63
against 0.65 --- the old device was a blunt proxy that happened to hit the right target,
and $\tau_C$ hits it directly. It also does what the old device structurally cannot:
express level uncertainty for 2006 and 2024, which is what limits the recent trend.

Numerically the two are near-equivalent for 1985, which is exactly why they must not be
stacked. On the 1985 campaign mean, the other terms alone give 0.0578; the old device
adds 0.0605 in quadrature and $\tau_C$ adds 0.0600. \textbf{Keeping both gives an
effective $\tau_C$ of 0.085} --- silently the 10\%/5\% broad scheme rather than the
6\%/3\% one adopted.

\begin{warnbox}[This run changes three things at once]
Correlation on, tails off, the 1985 inflation replaced by $\tau_C$. It is therefore a comparison of
\emph{error models}, not an isolation of the correlation effect, and the one-factor
discipline used in previous runs does not apply. Accepted deliberately: the
intermediate configurations are not scientifically interesting, and attribution can
be recovered afterwards from the conditional kurtosis and the campaign-specific
residual spreads.
\end{warnbox}

\section{Pre-registered predictions}

Stated before launch so that the outcome cannot be rationalised afterwards.

\begin{enumerate}[nosep]
\item $\sigma_{\text{input}}$ falls; MRT rises.
\item Posterior widths increase, most along the level direction.
\item RMSE is unchanged or slightly worse; the calibration--holdout gap narrows or
      holds.
\item $\sigma_{\text{init}}$ moves only weakly, and only via its correlation with
      $\sigma_{\text{input}}$.
\end{enumerate}

\textbf{Falsifier:} if posterior widths do \emph{not} increase, the implementation is
not doing what the theory says, whatever the medians do.

\begin{warnbox}[Comparison rules]
Log-likelihoods are \textbf{not} comparable across this change --- the normalising
constant moves. Compare on the RMSE distributions
(\texttt{S14\_rmse\_posterior}). And the standing rule applies with full force:
\textbf{$\tau_C$, $\tau_P$ and $\tau_R$ must never be tuned to land MRT on a
published value.}
\end{warnbox}

\section{What this is adopted for}

Not because it moves MRT. The independence assumption is measurably false: 1205
observations carry the weight of about 260 for the national level. The corrected
accounting is the more defensible error model whether or not the estimates move at
all. If MRT does not move, that is a result --- it would say the short residence
times are not an artefact of overconfidence, which is worth as much as the
alternative.

\end{document}
